% !TEX root =./ECE444_Practice_main.tex
%
% L11 practice set (High-Gain Antennas) -- aperture gain and beamwidth, reflector
% efficiency budgets and Ruze, Yagi-Uda behavior, and defending an antenna choice
% with a Friis link check.

\fullwidth{\textbf{Learning Objective 2.4: } I can describe how reflectors, Yagi-Uda
antennas, and arrays achieve high gain, and select an appropriate high-gain antenna
for a given application.
\\[4pt]\normalfont\small Show your work. A \textbf{Documentation} line at the end
is required (write \textbf{None} if you did not collaborate).}

\begin{questions}

%============================================================ Q1
\question \textbf{Sizing a C-band earth station.} A ground station uses a
\textbf{3.0\ \m\ diameter} prime-focus parabolic reflector at \textbf{6.0\ \ghz}.
Take the aperture efficiency to be $\eta_{\text{ap}} = 0.65$ and
$c = 3\ee{8}\ \text{m/s}$.

\begin{parts}

	\part In one or two sentences, explain why the statement ``gain is fundamentally a
	statement about area'' is more useful to a designer than ``gain is how much
	the antenna amplifies.''
		\begin{solution}[1.1in]
		An antenna is passive: it amplifies nothing. What it does is collect the
		radiated power into a smaller solid angle, and the size of the coherent
		radiating area --- measured in square wavelengths --- is what sets how
		small that solid angle can be. Writing $G = \eta_{\text{ap}} 4\pi A /
		\lambda^2$ tells the designer the only two knobs that exist: make the
		aperture physically bigger, or use a shorter wavelength. ``Amplification''
		suggests a third knob (more gain from the same box) that does not exist.
		\end{solution}

	\begin{minipage}[t]{0.58\linewidth}
		\part Find the wavelength, the dish diameter in wavelengths, and the
		gain in dBi.
		\begin{solution}[1.5in]
		\eq{
			\lambda &= \frac{c}{f} = \frac{3\ee{8}}{6\ee{9}} = 0.050\ \m
			\quad\Rightarrow\quad \frac{D}{\lambda} = \frac{3.0}{0.050} = 60 \\
			G &= \eta_{\text{ap}}\lp \frac{\pi D}{\lambda}\rp^{2}
			   = 0.65\,(60\pi)^{2} = 0.65\,(35531) \\
			  &= 2.31\ee{4}
			   \;\rightarrow\; 10\,\mylog{2.31\ee{4}} = 43.6\dbi
		}
		Remark: the effective aperture is $A_e = \eta_{\text{ap}}A =
		0.65\,\pi(1.5)^2 = 4.60\ \m^2$ --- that is the number Friis uses.
		\end{solution}
	\end{minipage}\hfill%
	\begin{minipage}[t]{0.37\linewidth}
		\raggedleft \vspace{12pt}
		$G$ = \ansbox[1.3in]{$43.6\dbi$}
	\end{minipage}

	\begin{minipage}[t]{0.58\linewidth}
		\part Estimate the half-power beamwidth, and state in one sentence what
		it implies about pointing.
		\begin{solution}[1.2in]
		\eq{
			\theta_{HP} \approx 70\mydeg\,\frac{\lambda}{D}
			= 70\mydeg\,\frac{0.050}{3.0} = 1.17\mydeg
		}
		A beam barely one degree wide means the mount has to hold the boresight
		to a small fraction of a degree; a few tenths of a degree of wind sway or
		thermal droop is already a measurable loss of signal. (Note also
		$2D^2/\lambda = 2(3)^2/0.05 = 360\ \m$ --- the far field starts a long way
		out, which is L12's problem.)
		\end{solution}
	\end{minipage}\hfill%
	\begin{minipage}[t]{0.37\linewidth}
		\raggedleft \vspace{12pt}
		$\theta_{HP}$ = \ansbox[1.3in]{$1.17\mydeg$}
	\end{minipage}

	\part Management wants \textbf{3\db\ more gain} at the same frequency. What
	dish diameter does that take, and name one specific cost of the change.
		\begin{solution}[1.3in]
		Gain goes as $D^2$, so $+3\db$ is a factor of 2 in gain and a factor of
		$\sqrt{2}$ in diameter:
		\eq{
			D_{new} = \sqrt{2}\,(3.0) = 4.24\ \m
			\qquad
			\theta_{HP,new} \approx 70\mydeg\,\frac{0.050}{4.24} = 0.83\mydeg
		}
		The cost: the beam narrows to $0.83\mydeg$, so the pointing budget
		tightens by the same factor of $\sqrt{2}$ --- and the reflector, mount,
		and wind load all grow with $D^2$. (Accept also: surface tolerance must
		be held over twice the area, or the far field moves out to 720\ \m.)
		\end{solution}

\end{parts}

\newpage
%============================================================ Q2
\question \textbf{Where the efficiency went.} Same 3.0\ \m\ dish. The
manufacturer quotes the following loss factors: spillover 0.88, illumination
taper 0.86, blockage 0.96, and a lumped ``everything else'' factor of 0.98. The
reflector surface is held to \textbf{0.5\ \text{mm}\ RMS}. Ruze's formula gives the
surface penalty as $\text{loss (dB)} = 685.8\,(\sigma/\lambda)^2$.

\begin{parts}

	\part Spillover and illumination taper trade directly against each other.
	Explain the trade in one or two sentences, and state the rule of thumb for
	where the optimum sits.
		\begin{solution}[1.2in]
		A narrow feed pattern puts almost all of its power on the dish, but
		leaves the rim many dB below the center, so the outer part of the
		aperture is not being used and the effective area shrinks. Widening the
		feed brightens the rim but sends power straight past it, which is lost
		entirely (and on receive replaces cold sky with 290\ K ground in the
		spilled beam). The optimum is the familiar rule of thumb: illuminate the
		\textbf{rim about 10\db\ below the center}.
		\end{solution}

	\begin{minipage}[t]{0.58\linewidth}
		\part Compute the aperture efficiency from the four quoted factors,
		before any surface error. Call it $\varepsilon_{ap,1}$.
		\begin{solution}[1.1in]
		\eq{
			\varepsilon_{ap,1} &= (0.88)(0.86)(0.96)(0.98) \\
			 &= 0.712
		}
		\end{solution}
	\end{minipage}\hfill%
	\begin{minipage}[t]{0.37\linewidth}
		\raggedleft \vspace{12pt}
		$\varepsilon_{ap,1}$ = \ansbox[1.3in]{$0.712$}
	\end{minipage}

	\begin{minipage}[t]{0.58\linewidth}
		\part Add the Ruze surface term at 6.0\ \ghz\ and report the total
		aperture efficiency.
		\begin{solution}[1.4in]
		At 6\ \ghz, $\lambda = 50\ \text{mm}$, so $\sigma/\lambda = 0.5/50 = 0.010$.
		\eq{
			\text{loss} &= 685.8\,(0.010)^{2} = 0.069\db \\
			\eta_{\text{surf}} &= 10^{-0.069/10} = 0.984 \\
			\eta_{\text{ap}} &= (0.712)(0.984) = 0.700
		}
		A half-millimetre surface is essentially free at C band.
		\end{solution}
	\end{minipage}\hfill%
	\begin{minipage}[t]{0.37\linewidth}
		\raggedleft \vspace{12pt}
		$\eta_{\text{ap}}$ = \ansbox[1.3in]{$0.70$}
	\end{minipage}

	\part Someone proposes reusing this \emph{same physical dish} at
	\textbf{30\ \ghz}. Recompute the surface term and the total aperture
	efficiency, then say in one or two sentences what the result means for the
	proposal.
		\begin{solution}[1.5in]
		At 30\ \ghz, $\lambda = 10\ \text{mm}$, so $\sigma/\lambda = 0.050$:
		\eq{
			\text{loss} &= 685.8\,(0.050)^{2} = 1.71\db \\
			\eta_{\text{surf}} &= 10^{-0.171} = 0.674 \\
			\eta_{\text{ap}} &= (0.712)(0.674) = 0.48
		}
		Nothing about the dish changed --- only $\lambda$ --- and the aperture
		efficiency fell from 0.70 to 0.48. The dish still gains about
		$56.3\dbi$ at 30\ \ghz\ (versus $58.0\dbi$ with a perfect surface), so
		the proposal is not absurd, but $1.7\db$ has been thrown away for free
		and it gets exponentially worse with frequency. Above roughly 30\ \ghz\
		this reflector is a machining problem, not an antenna.
		\end{solution}

\end{parts}

\newpage
%============================================================ Q3
\question \textbf{Reading a Yagi-Uda.} A Yagi-Uda has one driven element, one
reflector behind it, and a row of directors in front along a boom.

\begin{parts}

	\part The reflector is cut slightly \emph{longer} than resonance and the
	directors slightly \emph{shorter}. Explain what that detuning does
	electrically, and how it produces an endfire beam.
		\begin{solution}[1.4in]
		Only the driven element is connected; the parasites carry currents
		\emph{induced} by its near field, and the phase of that induced current
		is set by the element's reactance. A slightly long element is inductive,
		so its current lags; a slightly short element is capacitive, so its
		current leads. Placing the lagging element behind and the leading
		elements in front makes the re-radiated fields add in phase along the
		boom in the forward direction and tend to cancel in the reverse
		direction. The result is an endfire main lobe pointing away from the
		reflector, with a front-to-back ratio of roughly 15 to 25 dB.
		\end{solution}

	\part A second reflector buys almost nothing, but a seventh director still
	helps. Why?
		\begin{solution}[1.1in]
		The first reflector already cancels most of the field behind the
		antenna, so a second one sits in a region where there is almost no field
		to couple to and almost nothing left to cancel. The directors, by
		contrast, sit in the forward beam where the field is strong, and each one
		extends the slow forward travelling wave over a longer aperture. Gain
		follows the length of that aperture --- the boom --- not the number of
		rods on it.
		\end{solution}

	\begin{minipage}[t]{0.58\linewidth}
		\part A 6-element Yagi on a $1.0\lambda$ boom gives 10 dBi. Using the
		rule of thumb of $+3\db$ per doubling of boom length, estimate the boom
		needed for 16 dBi at 435\ \mhz. Then say whether you believe the
		answer.
		\begin{solution}[1.6in]
		$16 - 10 = 6\db$, which is two doublings, so the rule of thumb predicts
		a $4\lambda$ boom. At 435\ \mhz, $\lambda = 3\ee{8}/435\ee{6} = 0.69\ \m$:
		\eq{
			L = 4\lambda = 4\,(0.69) = 2.76\ \m
		}
		Do not believe it. The rule of thumb is optimistic past a few
		wavelengths: published designs put a $4.5\lambda$ boom at about
		14.5 dBi, not 16. Realistically you would either accept a much longer
		boom with diminishing returns, or stack two 13\dbi\ Yagis and take the
		$+3\db$ from the stack.
		\end{solution}
	\end{minipage}\hfill%
	\begin{minipage}[t]{0.37\linewidth}
		\raggedleft \vspace{12pt}
		$L$ = \ansbox[1.3in]{$2.76\ \m$}
	\end{minipage}

	\part Yagis are narrowband --- a few percent. Explain why, and name one
	application where that does not matter.
		\begin{solution}[1.1in]
		Every element is a detuned resonator, and the whole design depends on the
		\emph{phases} of the induced parasite currents. Move off the design
		frequency and each element's reactance changes, the phases drift, and the
		forward addition and rearward cancellation both degrade --- gain drops and
		the front-to-back ratio collapses. That is fine anywhere the frequency is
		fixed: a single-channel broadcast TV antenna, an amateur band antenna, or
		a fixed point-to-point link.
		\end{solution}

\end{parts}

\newpage
%============================================================ Q4
\question \textbf{Choosing an antenna.} You need a \textbf{25\dbi} antenna at
each end of a 2\ \km\ rooftop-to-rooftop link at \textbf{5.8\ \ghz}. The mounts
are fixed (no steering), the site is windy, and the budget is small.

\begin{parts}

	\begin{minipage}[t]{0.58\linewidth}
		\part Find the required effective aperture, then the diameter of a
		parabolic dish that delivers it with $\eta_{\text{ap}} = 0.60$.
		\begin{solution}[1.6in]
		$\lambda = 3\ee{8}/5.8\ee{9} = 0.0517\ \m$ and $G = 10^{25/10} = 316$.
		\eq{
			A_e &= \frac{G\lambda^{2}}{4\pi}
			     = \frac{316\,(0.0517)^{2}}{4\pi} = 0.0673\ \m^{2} \\
			A   &= \frac{A_e}{\eta_{\text{ap}}} = \frac{0.0673}{0.60}
			     = 0.112\ \m^{2} \\
			D   &= 2\sqrt{A/\pi} = 2\sqrt{0.0357} = 0.378\ \m
		}
		Beamwidth: $\theta_{HP} \approx 70\mydeg(0.0517/0.378) = 9.6\mydeg$.
		\end{solution}
	\end{minipage}\hfill%
	\begin{minipage}[t]{0.37\linewidth}
		\raggedleft \vspace{12pt}
		$D$ = \ansbox[1.3in]{$0.38\ \m$}
	\end{minipage}

	\begin{minipage}[t]{0.58\linewidth}
		\part A vendor offers a flat patch array instead, with
		$\eta_{\text{ap}} = 0.75$ and elements on a $\lambda/2$ grid. How many
		elements, and how big is the panel?
		\begin{solution}[1.5in]
		\eq{
			A &= \frac{A_e}{0.75} = \frac{0.0673}{0.75} = 0.0898\ \m^{2}
			\quad\Rightarrow\quad \text{side} = 0.300\ \m \\
			\frac{\lambda}{2} &= 0.0259\ \m
			\quad\Rightarrow\quad \frac{0.300}{0.0259} = 11.6 \;\rightarrow\; 12
		}
		A $12\times12 = 144$ element panel, about $0.31\ \m$ square. Check:
		$G = 0.75\,(4\pi)(0.31)^2/(0.0517)^2 = 339 = 25.3\dbi$. Good.
		\end{solution}
	\end{minipage}\hfill%
	\begin{minipage}[t]{0.37\linewidth}
		\raggedleft \vspace{12pt}
		$N$ = \ansbox[1.3in]{$144$}
	\end{minipage}

	\part Pick one and defend it in two or three sentences, using the numbers
	you just computed.
		\begin{solution}[1.4in]
		Take the dish. It is a single moulded part 0.38\ \m\ across with one
		feed, it is broadband, and its wind load at that size is modest; the
		patch array needs 144 elements and a 144-way corporate feed network whose
		loss at 5.8\ \ghz\ is significant in both cost and dB. The array's
		advantages --- flat profile and the option of electronic steering --- are
		worth nothing
		on a fixed rooftop link, so they do not justify the cost. (A defensible
		opposite answer: choose the panel if the roof lease or the wind survey
		forbids a dish, and accept the feed loss.)
		\end{solution}

	\begin{minipage}[t]{0.58\linewidth}
		\part The dish beam is $9.6\mydeg$ wide. Using the beam-shape
		approximation $L(\theta) \approx 3\,(2\theta/\theta_{HP})^{2}$ dB, how
		far off boresight can the mount drift before you lose 0.5 dB?
		\begin{solution}[1.2in]
		\eq{
			0.5 &= 3\lp\frac{2\theta}{9.6\mydeg}\rp^{2}
			\quad\Rightarrow\quad
			\frac{2\theta}{9.6\mydeg} = \sqrt{0.1667} = 0.408 \\
			\theta &= \frac{(0.408)(9.6\mydeg)}{2} = 1.96\mydeg
		}
		About $\pm2\mydeg$ --- comfortable for a rigid rooftop mount, which is
		another argument for not buying more gain than the link needs.
		\end{solution}
	\end{minipage}\hfill%
	\begin{minipage}[t]{0.37\linewidth}
		\raggedleft \vspace{12pt}
		$\theta$ = \ansbox[1.3in]{$\pm 2.0\mydeg$}
	\end{minipage}

\end{parts}

\newpage
%============================================================ Q5
\question \textbf{Does the link close?} Continue the 5.8\ \ghz, 2\ \km\ link
from the previous question. The transmitter delivers \textbf{20 dBm} and the
receiver sensitivity is \textbf{$-80$ dBm}. Use Friis,
$P_r = P_t G_t G_r \lp\lambda/4\pi R\rp^{2}$.

\begin{parts}

	\begin{minipage}[t]{0.58\linewidth}
		\part Compute the free-space path loss in dB.
		\begin{solution}[1.3in]
		\eq{
			L_{fs} &= 20\,\mylog{\frac{4\pi R}{\lambda}}
			 = 20\,\mylog{\frac{4\pi(2000)}{0.0517}} \\
			 &= 20\,\mylog{4.86\ee{5}} = 113.7\db
		}
		\end{solution}
	\end{minipage}\hfill%
	\begin{minipage}[t]{0.37\linewidth}
		\raggedleft \vspace{12pt}
		$L_{fs}$ = \ansbox[1.3in]{$113.7\db$}
	\end{minipage}

	\begin{minipage}[t]{0.58\linewidth}
		\part Find the received power with the 25\dbi\ dishes at both ends, and
		the link margin.
		\begin{solution}[1.3in]
		\eq{
			P_r &= P_t + G_t + G_r - L_{fs} \\
			 &= 20 + 25 + 25 - 113.7 = -43.7\dbm \\
			\text{margin} &= -43.7 - (-80) = 36.3\db
		}
		\end{solution}
	\end{minipage}\hfill%
	\begin{minipage}[t]{0.37\linewidth}
		\raggedleft \vspace{12pt}
		$P_r$ = \ansbox[1.3in]{$-43.7\dbm$}
	\end{minipage}

	\begin{minipage}[t]{0.58\linewidth}
		\part Now replace both dishes with 8\dbi\ patches. Recompute $P_r$ and
		the margin.
		\begin{solution}[1.2in]
		\eq{
			P_r &= 20 + 8 + 8 - 113.7 = -77.7\dbm \\
			\text{margin} &= -77.7 - (-80) = 2.3\db
		}
		The 34\db\ of antenna gain you gave up came straight out of the margin.
		\end{solution}
	\end{minipage}\hfill%
	\begin{minipage}[t]{0.37\linewidth}
		\raggedleft \vspace{12pt}
		$P_r$ = \ansbox[1.3in]{$-77.7\dbm$}
	\end{minipage}

	\part Which link would you actually deploy, and why? Name two specific
	things that eat margin on a real rooftop link.
		\begin{solution}[1.3in]
		Deploy the dish link. A 2.3\db\ margin leaves no room for normal
		propagation variation, so the link fails on the first bad day. Real
		margin eaters include rain and wet foliage in the path, mount drift and building sway pushing the beam off
		boresight (Q4d showed only $\pm2\mydeg$ of slack), multipath from nearby
		rooftops, and connector and cable loss between the radio and the
		antenna. With 36\db\ in hand the dish link absorbs all of that; with
		2.3\db\ it absorbs none of it.
		\end{solution}

\end{parts}

\vspace{1.5em}
\noindent\textbf{Documentation:} \rule[-2pt]{0.6\linewidth}{0.4pt}

\end{questions}
